\documentclass[10pt]{amsart}
\usepackage[a4paper,margin=22mm]{geometry}
\usepackage[no-math]{fontspec}
\setmainfont{Latin Modern Roman}[SmallCapsFont={lmromancaps10-regular.otf}]
\usepackage{amsmath,amssymb,amsthm,mathtools,hyperref,graphicx}
\usepackage[shortlabels]{enumitem}
\emergencystretch=3em
\newtheorem{theo}{Theorem}[section]
\newtheorem{lemm}[theo]{Lemma}
\theoremstyle{definition}
\newtheorem{defi}[theo]{Definition}
\newtheorem{rema}[theo]{Remark}
% Published pagination breaks do not control the independent standalone layout.
\renewcommand{\pagebreak}[1][4]{}
\newcommand{\E}{\mathbb{E}}
\newcommand{\G}{\mathbb{G}}
\newcommand{\V}{\mathbb{V}}
\newcommand{\Z}{\mathbb{Z}}
\newcommand{\Pro}{\mathbb{P}}


\newcommand{\clF}{\mathcal{F}}
\newcommand{\ee}{\mathbf{e}}
\newcommand{\calC}{\mathcal{C}}
\newcommand{\calG}{\mathcal{G}}
\newcommand{\Expl}{\mathrm{Expl}}
\newcommand{\diam}{\mathrm{diam}}






%%%%%%%%%%%%%%%%%%%%%%%%%%%%%%%%%%%%%%%%%%%%%%%%%%%%%%%%%%%%%%%%%%%%%%%%%%%%%
\graphicspath{{./figures/}}

\newcommand*{\mk}{\mkern -1mu}
\newcommand*{\Mk}{\mkern -2mu}
\newcommand*{\mK}{\mkern 1mu}
\newcommand*{\MK}{\mkern 2mu}

\hypersetup{urlcolor=purple, linkcolor=blue, citecolor=red}

\let\oldforall\forall
\renewcommand*{\forall}{\mathrel{\oldforall}}
\title[Ghost-field near-critical sharpness]{Ghost-field near-critical sharpness bound for Bernoulli percolation on connected locally finite vertex-transitive countably infinite graphs}
\author{Hugo Vanneuville}
\date{Annales Henri Lebesgue 8 (2025),101–112; DOI 10.5802/ahl.230}
\hypersetup{hidelinks}
\begin{document}
\maketitle
\section{Introduction}

Let us consider \emph{Bernoulli bond percolation} on a connected, locally finite, vertex-transitive and countably infinite graph $\G=(\V,\E)$ (e.g.\ the hypercubic lattice $\Z^d$). ``Locally finite'' means that the degree of each vertex is finite, and ``vertex-transitive'' means that for all vertices $v,w$, there exists an automorphism of $\G$ that maps~$v$ to~$w$.


This model is defined as follows: each edge is declared \emph{open} with some probability $p$ and \emph{closed} otherwise, independently of the other edges. We let $\Pro_p$ denote the corresponding product probability measure on $\{0,1\}^{\E}$, where $1$ means open and $0$ means closed. In percolation theory, one is interested in the connectivity properties of the graph obtained by keeping only the open edges. We fix a vertex $o \in \V$ once and for all, we let $\calC_o$ denote the cluster of $o$, i.e. the set of all vertices that are connected to $o$ by an open path, and we let $|\calC_o|$ denote its cardinality. The \emph{critical probability} is defined as follows:
\begin{align*}
p_c &= \inf \big\{ p \in [0,1] : \Pro_p\big[|\calC_o|=+\infty\big]>0 \big\}.
\\
\intertext{Let}
\psi_n(p) &= \Pro_p\big[|\calC_o| \ge n\big], \quad \forall n \ge 0.
\end{align*}
To study the volume of clusters, it is relevant to introduce what is often referred to as a ``ghost field''. By this, we just mean that, given some parameter $h \in (0,+\infty)$, we color every vertex \emph{green} with probability $1-e^{-h}$, independently of the other vertices and of the bond percolation configuration. We let $\Pro_{p,h}$ denote the product probability measure on $\{0,1\}^{\E} \times \{0,1\}^{\V}$ with this additional feature, where we assign the number $1$ to every green vertex, i.e.
\[
\Pro_{p,h} = \big(p\delta_1+(1-p)\delta_0\big)^{\otimes \E} \otimes \big((1-e^{-h})\delta_1+e^{-h}\delta_0\big)^{\otimes \V}.
\]
We let $\calG$ denote the set of green vertices and we consider the so-called \emph{magnetization} (this terminology comes from a comparison with analogous quantities that appear in the Ising model):
\[
m_h(p)=\Pro_{p,h} \big[\calC_o \cap \calG \ne \emptyset \big].
\]
One can note that, for every $p$, $m_h(p)$ goes to $\Pro_p[|\calC_o|=+\infty]$ as $h$ goes to $0$.
\setcounter{theo}{1}
\begin{theo}\label{thm:exp}
Let $p\in(0,1)$ and $h\in(0,+\infty)$, and let $q=p(1-m_h(p))$. Then,
\[
\forall n \ge 0, \quad \psi_n(q) \le \frac{1}{1-m_h(p)} \psi_n(p) e^{-hn}.
\]
\end{theo}
\setcounter{section}{2}
\section{A general criterion to compare the effect of decreasing \texorpdfstring{$p$}{p} to that of conditioning on an event}\label{sec:gen}

\subsection{Explorations}

Let us fix a finite graph $G=(V,E)$ and call $\Pro_p^G$ and $\Pro_{p,h}^{G}$ the analogues of $\Pro_p$ and $\Pro_{p,h}$ on the graph $G$.\footnote{I.e.\ $\Pro_p^G$ is the product Bernoulli measure of parameter $p$ on $\{0,1\}^E$ and $\Pro_{p,h}^G$ is the product Bernoulli measure on $\{0,1\}^E\times\{0,1\}^V$, of parameter $p$ on the edges and $1-e^{-h}$ on the vertices. The proofs in this section actually work in a more general setting, for instance by considering a probability space $(M,\clF,\rho)$ and considering the measure $\Pro_p^G\otimes\rho$ instead of $\Pro_{p,h}^G$.} We denote by $\omega$ the elements of $\{0,1\}^E$ and by $\eta$ the elements of $\{0,1\}^V$. Moreover, we let $\vec{E}$ be the set of all orderings $(e_1,\dots,e_{|E|})$ of $E$.

\begin{defi}\label{defi:expl}
An \emph{exploration} of $E$ is a map
\begin{align*}
\ee : \{0,1\}^E & \longrightarrow \vec{E}\\
\omega & \longmapsto (\ee_1,\dots,\ee_{|E|})=\big(\ee_1(\omega),\dots,\ee_{|E|}(\omega)\big)
\end{align*}
such that $\ee_1$ does not depend on $\omega$ and $\ee_{k+1}$ only depends on $(\ee_1,\dots,\ee_k)$ and $(\omega_{\ee_1},\dots,\omega_{\ee_k})$.\footnote{This means that for every $k \in \{0,\dots,|E|-1\}$, there exists a function $\phi_k$ such that $\forall \omega \in \{0,1\}^E, \; \ee_{k+1}(\omega)=\phi_k\big(\ee_1(\omega),\dots,\ee_k(\omega),\omega_{\ee_1(\omega)},\dots,\omega_{\ee_k(\omega)})$.}
Given an exploration $\ee$, $(e,x) \in \vec{E}\times\{0,1\}^E$ and $k \in \{0,\dots,|E|\}$, we denote by $\Expl_k(e,x)$ the event that $(\ee,\omega)$ coincides with $(e,x)$ at least until step $k$, i.e.
\[
\Expl_k(e,x)=\big\{\omega \in \{0,1\}^E : \forall j \in \{1,\dots,k\}, \ee_j=e_j\text{ and }\omega_{e_j}=x_{e_j}\big\}.
\]
(In particular, $\Expl_0(e,x)=\{0,1\}^E$.) We also let $\Expl(e,x)=\Expl_{|E|}(e,x)$.
\end{defi}

We recall that an \emph{increasing} subset $A\subseteq\{0,1\}^E$ is a subset that satisfies ($\omega \in A$ and $\forall e \in E, \quad \omega'_e \ge \omega_e$) $\Rightarrow$ $\omega'\in A$. Moreover, given two probability measures $\mu,\nu$ on $\{0,1\}^E$, we say that $\mu$ is \emph{stochastically smaller} than $\nu$ if $\mu[A]\le \nu[A]$ for every increasing subset $A$, and we denote this property by $\mu \preceq \nu$. Finally, we say that an edge $e \in E$ is \emph{pivotal} for a set $A \subseteq \{0,1\}^E \times \{0,1\}^V$ and an element~$(\omega,\eta)$ if changing only the edge $e$ in $\omega$ (and not changing $\eta$) modifies~$1_A(\omega,\eta)$.

\subsection{The main intermediate lemma}

The key result of this section is the following lemma. We fix an exploration of $E$ for the rest of the section.

\begin{lemm}\label{lem:coupl}
Let $p \in (0,1)$ and $h \in (0,+\infty)$. Moreover, let $A \subseteq \{0,1\}^E \times \{0,1\}^V$ be any non-empty set and let $\varepsilon \in [0,1]$. Assume that for every $(e,x) \in \vec{E}\times\{0,1\}^E$ such that $A \cap \Expl(e,x)\neq \emptyset$, we have
\[
\forall k \in \{0,\dots,|E|-1\}, \quad \Pro_{p,h}^G \big[\text{$e_{k+1}$ is pivotal for $A$} \; \big| \; A \cap \Expl_k(e,x) \big] \le \varepsilon.
\]
Then,
\[
\Pro_{p(1-\varepsilon)}^G \preceq \Pro_{p,h}^G \big[\omega \in \cdot \; \big| \; A \big].
\]
\end{lemm}

\begin{rema}\label{rk:formal}\leavevmode
\begin{itemize}
\item Recall that we use the letter $\omega$ to denote the elements of $\{0,1\}^E$. So $\Pro_{p,h}^G [\omega \in \cdot \mid A]$ denotes the marginal law on $\{0,1\}^E$ of $\Pro_{p,h}^G[\, \cdot \mid A]$;
\item We have written $\Expl_k(e,x)$ when we should have written $\Expl_k(e,x) \times \{0,1\}^V$;
\item Let us take the opportunity of this remark to make the following observation that holds when $(e,x)$ is in the image of $\omega \mapsto (\ee(\omega),\omega)$. In this case, $\Expl_k(e,x)=\{\omega \in \{0,1\}^E : \forall j \in \{1,\dots,k\}, \omega_{e_j}=x_{e_j}\}$. As a result, $\Pro_{p,h}^G [\, \cdot \mid \Expl_k(e,x)]$ is the probability measure on $\{0,1\}^E\times\{0,1\}^V$ that assigns the value $x_{e_j}$ to $e_j$ for every $j \le k$ and is still the product Bernoulli measure of parameter $p$ on the other edges and of parameter $1-e^{-h}$ on the vertices.
\end{itemize}
\end{rema}

The proof of Lemma~\ref{lem:coupl} is based on the following lemma (see e.g.~\cite[Lemma~1]{Rus82}, \cite[(7.64)]{Gri99} and~\cite[Lemma~2.1]{DRT19} for very similar results).

\begin{lemm}\label{lem:gen}
Let $q \in [0,1]$. Moreover, let $\mu$ be a probability measure on $\{0,1\}^E$ and assume that for every $(e,x) \in \vec{E}\times\{0,1\}^E$ such that $\mu[\Expl(e,x)]>0$, we have
\[
\forall k \in \{0,\dots,|E|-1\}, \quad \mu \big[\omega_{e_{k+1}} = 1 \; \big| \; \Expl_k(e,x) \big] \ge q.
\]
Then, $\Pro_q^G \preceq \mu$.
\end{lemm}

\begin{proof}[Proof of Lemma~\ref{lem:gen}]
The result can be proven by induction on $|E|$. Throughout the proof, it is important to remember that $\ee_1$ is constant. Let $A\subseteq\{0,1\}^E$ be an increasing set. We want to prove that $\mu[A]\ge\Pro^G_q[A]$. We extend every $\omega \in \{0,1\}^{E\setminus\{\ee_1\}}$ into two configurations $\omega^0,\omega^1\in\{0,1\}^E$ by setting $\omega^i_{\ee_1}=i$. We define $A^0,A^1\subseteq\{0,1\}^{E\setminus\{\ee_1\}}$ by $A^i=\{\omega \in \{0,1\}^{E\setminus\{\ee_1\}} : \omega^i \in A\}$. Moreover, we define two probability measures $\mu^0,\mu^1$ on $\{0,1\}^{E\setminus\{\ee_1\}}$ by $\mu^i=\mu[\omega_{|E\setminus\{\ee_1\}}\in\cdot\mid\omega_{\ee_1}=i]$.


Let $G'$ be the graph obtained from $G$ by erasing $\ee_1$. By the induction hypothesis -- applied to the explorations $\omega\in\{0,1\}^{E\setminus\{\ee_1\}}\mapsto(\ee_2(\omega^i),\dots,\ee_{|E|}(\omega^i))$ --, we have $\mu^i[A^i]\ge \Pro^{G'}_q[A^i]$ for every $i\in \{0,1\}$. As a result,
\begin{align*}
\mu[A] &=\mu[\omega_{\ee_1}=0]\mu^0\big[A^0\big]+ \mu[\omega_{\ee_1}=1]\mu^1\big[A^1\big]\\
&\ge \mu[\omega_{\ee_1}=0]\Pro_q^{G'} \big[A^0\big]+ \mu[\omega_{\ee_1}=1]\Pro_q^{G'} \big[A^1\big].
\end{align*}
By the assumption of the lemma for $k=0$, we have $ \mu[\omega_{\ee_1}=1] \ge q$. Moreover, $A^0\subseteq A^1$ since $A$ is increasing, so the above is larger than or equal to
\[
(1-q)\Pro_q^{G'} \big[A^0\big]+q \Pro_q^{G'} \big[A^1\big] = \Pro^G_q[A]. \qedhere
\]
\end{proof}

\begin{proof}[Proof of Lemma~\ref{lem:coupl}]
Let us prove that the hypothesis of Lemma~\ref{lem:gen} holds with $\mu=\Pro_{p,h}^G[\omega \in \cdot \mid A]$ and $q=p(1-\varepsilon)$. Let $(e,x)$ as in the statement of Lemma~\ref{lem:coupl}. The third point of Remark~\ref{rk:formal} implies that, under $\Pro_{p,h}^G [\, \cdot \mid \Expl_k(e,x) ]$, $\omega_{e_{k+1}}$ is independent of the event $A \cap \{ \text{$e_{k+1} $ is not piv. for $A$} \}$ (because the latter depends only on $\eta$ and $\omega_{|E\setminus\{e_{k+1}\}}$). As a result,
\pagebreak
\begin{multline*}
\Pro_{p,h}^G \big[\omega_{e_{k+1}} = 1 \; \big| \; A \cap \Expl_k(e,x) \big]\\
\ge \frac{\Pro_{p,h}^G \big[\omega_{e_{k+1}} = 1, A, \text{$e_{k+1} $ is not piv. for $A$} \; \big| \; \Expl_k(e,x) \big]}{\Pro_{p,h}^G \big[A \; \big| \; \Expl_k(e,x) \big]}\\
= p \times\Pro_{p,h}^G \big[\text{$e_{k+1} $ is not piv. for $A$} \; \big| \; A \cap \Expl_k(e,x) \big]\ge p~(1-\varepsilon).
\end{multline*}
We end the proof by applying Lemma~\ref{lem:gen}.
\end{proof}

\section{Proof of Theorem~\ref{thm:exp}}\label{sec:proof}

Let $G_n=(V_n,E_n)$ denote the graph $\G$ restricted to the ball of radius $n$ around~$o$ (for the graph distance). We keep the notations~$\calC_o$ and~$\calG$ (the cluster of $o$ and the ghost field) in the context of percolation on $G_n$. As in Section~\ref{sec:gen}, $\Pro_{p}^{G_n}$ and $\Pro_{p,h}^{G_n}$ are the analogues of $\Pro_p$ and $\Pro_{p,h}$ on the graph $G_n$, and we denote by $\omega$ the elements of $\{0,1\}^{E_n}$.

\begin{lemm}\label{lem:coupl_perco}
Let $p \in (0,1)$ and $h \in (0,+\infty)$, and write $q=p(1-m_h(p))$. Then,
\[
\forall n \ge 0, \quad \Pro_{q}^{G_n} \preceq \Pro_{p,h}^{G_n} \big[\omega \in \cdot \; \big| \; \calC_o \cap \calG = \emptyset \big].
\]
\end{lemm}

\begin{proof}
Let us apply Lemma~\ref{lem:coupl}, with $A=\{\calC_o\cap\calG=\emptyset\}$ and $\varepsilon=m_h(p)$. To this purpose, we define an exploration of $E_n$ by fixing an arbitrary ordering of $E_n$ and:
\begin{itemize}
\item By first revealing $\calC_o$ (i.e.\ by revealing iteratively the edge of smallest index among all the edges that are connected to $o$ by already revealed open edges). Here and below, `smallest index' refers to the arbitrary ordering that we have fixed;
\item Then, by revealing all the other edges (once again by revealing iteratively the edge of smallest index).
\end{itemize}
Let $(e,x) \in \vec{E}_n\times\{0,1\}^{E_n}$ such that $\{\calC_o\cap\calG=\emptyset\}\cap\Expl(e,x)\neq\emptyset$. We observe that, if $\Expl_k(e,x)$ holds and $e_{k+1}$ is pivotal for $\{\calC_o \cap \calG = \emptyset\}$, then one can write $e_{k+1}=\{v_{k+1},w_{k+1}\}$ where:
\begin{itemize}
\item $v_{k+1}$ is connected to $o$ by open edges in $\{e_1,\dots,e_k\}$ but $w_{k+1}$ is not;
\item $w_{k+1}$ is connected to a green vertex by open edges that do not belong to $\{e_1,\dots,e_{k+1}\}$.
\end{itemize}
We let $B_{k+1}$ denote the event of the second item. By this observation and then the Harris--FKG inequality\footnote{If one does not want to use the Harris--FKG inequality in this paper, one can continue to explore the configuration with the same rule as before, except that $e_{k+1}$ is not revealed until no other unrevealed edge is connected to $o$, and one can compute the pivotal probability conditionally on this further information.} (see for instance~\cite[Theorem~2.4]{Gri99} or~\cite[Lemma~3]{BR06}) applied to $\Pro_{p,h}^{G_n} \big[\, \cdot \mid \Expl_k(e,x)\big]$ (which is a product of Bernoulli laws by the third point of Remark~\ref{rk:formal}) and to the increasing event $B_{k+1}$ and the decreasing event $\{\calC_o\cap\calG=\emptyset\}$, we obtain that
\pagebreak
\begin{multline*}
\Pro_{p,h}^{G_n} \big[e_{k+1}\text{ is piv.\ for } \{\calC_o \cap \calG = \emptyset\} \; \big| \; \calC_o \cap \calG = \emptyset, \Expl_k(e,x) \big]\\
\begin{aligned}
&\le \Pro_{p,h}^{G_n} \big[B_{k+1} \; \big| \; \calC_o \cap \calG = \emptyset, \Expl_k(e,x) \big]\\
&\mkern-23mu\overset{\text{Harris--FKG}}{\le} \Pro_{p,h}^{G_n} \big[B_{k+1} \; \big| \; \Expl_k(e,x) \big]= \Pro_{p,h}^{G_n} \big[B_{k+1} \big] \le m_h(p),
\end{aligned}
\end{multline*}
where in the equality we have used the third point of Remark~\ref{rk:formal}. We conclude by applying Lemma~\ref{lem:coupl}.
\end{proof}

\begin{proof}[Proof of Theorem~\ref{thm:exp}]
We note that $\psi_n(q)=\Pro_q^{G_n} \big[|\calC_o|\ge n \big]$. As a result, Lemma~\ref{lem:coupl_perco} implies that
\begin{align*}
\psi_n(q)\le\Pro_{p,h}^{G_n} \big[|\calC_o|\ge n \; \big| \; \calC_o \cap \calG = \emptyset
\big] &\overset{\text{Bayes}}{=} \frac{\psi_n(p) \times \Pro_{p,h}^{G_n}\big[\calC_o \cap \calG = \emptyset \; \big| \; |\calC_o| \ge n \big]}{1-\Pro_{p,h}^{G_n}\big[\calC_o \cap \calG \neq \emptyset \big]}\\
& \hspace{0.2cm}= \frac{\psi_n(p) \times \E_{p,h}^{G_n}\big[e^{-h|\calC_o|} \; \big| \; |\calC_o| \ge n \big]}{1-\Pro_{p,h}^{G_n}\big[\calC_o \cap \calG \neq \emptyset \big]}\\
& \hspace{0.2cm}\le \frac{\psi_n(p)}{1-m_h(p)} e^{-hn}. \qedhere
\end{align*}
\end{proof}
\begin{thebibliography}{DCRT19}
\bibitem[BR06]{BR06} Bollobás, Béla; Riordan, Oliver Percolation, Cambridge University Press, 2006.
\bibitem[DCRT19]{DRT19} Duminil-Copin, Hugo; Raoufi, Aran; Tassion, Vincent Sharp phase transition for the random-cluster and Potts models via decision trees, Ann. Math. (2), Volume 189 (2019) no. 1, pp. 75-99.
\bibitem[Gri99]{Gri99} Geoffrey R. Grimmett, Percolation, 2nd ed., Grundlehren der Mathematischen Wissenschaften, vol.321, Springer, 1999.
\bibitem[Rus82]{Rus82} Russo, Lucio An approximate zero-one law, Z. Wahrscheinlichkeitstheor. Verw. Geb., Volume 61 (1982), pp. 129-139.
\end{thebibliography}
\end{document}
